\chapter{Effective descent: restriction of the glued sheaf to a chart}
\label{chap:Picard_GlueDescent}

\providecommand{\ShMod}{\mathrm{SheafOfModules}}

The gluing primitive \cref{def:scheme_modules_glue} produces, out of a glue datum
\(D\) of schemes (index set \(J\), charts \(U_i\), overlaps \(V_{ij}\), overlap
immersions \(f_{ij} : V_{ij} \hookrightarrow U_i\), transition isomorphisms
\(t_{ij} : V_{ij} \xrightarrow{\sim} V_{ji}\), glued scheme \(X\) with chart
immersions \(\iota_i : U_i \hookrightarrow X\)), a sheaf of \(\mathcal{O}_X\)-modules
\(\mathcal{M}\) as the descent equalizer
\[
  \mathcal{M} \;=\; \operatorname{eq}(a, b)
    \;\hookrightarrow\; \prod_{i} (\iota_i)_{*}\,\mathcal{M}_i
\]
of two parallel maps into the overlap product. That construction makes
\(\mathcal{M}\) but does \emph{not} yet consume the cocycle conditions (C1)/(C2) on
the transition family \((g_{ij})\). This chapter discharges the remaining content of
descent: it shows that restricting the glued sheaf back to a chart recovers the
input sheaf there. Concretely, the canonical \emph{restriction morphism}
\(\iota_i^{*}\mathcal{M} \to \mathcal{M}_i\) -- the adjoint transpose along the chart
immersion \(\iota_i\) of the \(i\)-th equalizer projection -- is an isomorphism, and
it is here that (C1)/(C2) are used. The engine is that pullback along an open
immersion is, up to natural isomorphism, the site-level restriction functor
(\cref{lem:modules_restrictFunctorIsoPullback_mathlib}); being a right adjoint it
preserves the descent equalizer and the pushforward product, so the chart
restriction of \(\mathcal{M}\) is computed factor by factor. The non-formal input is
the \emph{overlap base change} \(\beta_{ij}\), which identifies the restricted
pushforward factor \(\iota_i^{*}(\iota_j)_{*}\mathcal{M}_j\) with a pushforward of a
restriction of \(\mathcal{M}_j\) to the overlap, using the cartesianness of the
chart--overlap square.

\section{The descent equalizer and its overlap base change}
\label{sec:gd_equalizer}

\begin{definition}

\leanok
[The pushforward product carrying the descent equalizer]
  \label{def:gr_glue_equalizer}
  \lean{AlgebraicGeometry.Scheme.Modules.glueProd}
  \uses{def:scheme_modules_glue}
  Let \(D\) be a glue datum and \((\mathcal{M}_i)_{i \in J}\) a family of sheaves of
  modules with \(\mathcal{M}_i\) on the chart \(U_i\). Write
  \[
    P \;=\; \prod_{i} (\iota_i)_{*}\,\mathcal{M}_i
    \qquad\text{and}\qquad
    Q \;=\; \prod_{(i,j)} (f_{ij} \circ \iota_i)_{*}\,\bigl(f_{ij}^{*}\mathcal{M}_i\bigr)
  \]
  for the \emph{pushforward product} and the \emph{overlap product} of sheaves of
  \(\mathcal{O}_X\)-modules. The two descent legs of \cref{def:scheme_modules_glue},
  \[
    P \;\underset{b}{\overset{a}{\rightrightarrows}}\; Q,
  \]
  are: on the \((i,j)\)-component, leg \(a\) restricts the \(i\)-th chart factor to
  the overlap \(V_{ij}\) using the unit of the pullback--pushforward adjunction along
  \(f_{ij}\) and the pushforward-composition comparison
  \((\iota_i)_{*} \cong (f_{ij} \circ \iota_i)_{*} \circ f_{ij}^{*}\); leg \(b\)
  restricts the \(j\)-th chart factor, transports it across the inverse transition
  isomorphism \(g_{ij}^{-1}\), and reindexes the immersion via the glue condition
  \((t_{ij} \circ f_{ji}) \circ \iota_j = f_{ij} \circ \iota_i\). The glued sheaf
  \(\mathcal{M} = \operatorname{eq}(a,b)\) is the equalizer of this parallel pair;
  \(P\) is the object it embeds into and \(Q\) the object detecting the overlap
  agreement. The product \(P\) is the carrier on which the \(i\)-th projection (and,
  through it, the restriction morphism) is defined.
\end{definition}

\begin{lemma}

\leanok
[Overlap base change of the chart square (\(\beta_{ij}\))]
  \label{lem:glueOverlapBaseChangeIso}
  \lean{AlgebraicGeometry.Scheme.Modules.glueOverlapBaseChangeIso}
  \uses{lem:modules_restrictFunctorIsoPullback_mathlib}
  For chart indices \(i, j \in J\) there is a natural isomorphism of functors
  \(\ShMod(\mathcal{O}_{U_j}) \to \ShMod(\mathcal{O}_{U_i})\),
  \[
    \beta_{ij} \;:\;
      \iota_i^{*}\,(\iota_j)_{*}
        \;\xrightarrow{\ \sim\ }\;
      (f_{ij})_{*}\,(t_{ij} \circ f_{ji})^{*},
  \]
  identifying, for any sheaf \(\mathcal{N}\) on the chart \(U_j\), the restriction to
  the chart \(U_i\) of the pushforward \((\iota_j)_{*}\mathcal{N}\) with the
  pushforward along \(f_{ij}\) of the restriction of \(\mathcal{N}\) to the overlap
  \(V_{ij}\) along \(t_{ij} \circ f_{ji}\). Here the chart restriction \(\iota_i^{*}\)
  and the overlap restriction \((t_{ij} \circ f_{ji})^{*}\) are realised as the
  site-level restriction functors, isomorphic to the geometric pullbacks by
  \cref{lem:modules_restrictFunctorIsoPullback_mathlib}.
\end{lemma}
\begin{proof}
  \uses{lem:modules_restrictFunctorIsoPullback_mathlib, lem:gr_overlap_appIso_compat,
    lem:gr_overlapBaseChange_inv_app, lem:gr_overlapBaseChange_hom_app}
    \leanok
  The isomorphism is the cartesianness of the chart--overlap square, made site-level.
  The square
  \[
    \begin{array}{ccc}
      V_{ij} & \xrightarrow{\ t_{ij} \circ f_{ji}\ } & U_j \\
      \ \downarrow{\scriptstyle f_{ij}} & & \downarrow{\scriptstyle \iota_j} \\
      U_i & \xrightarrow{\ \ \iota_i\ \ } & X
    \end{array}
  \]
  is cartesian, and its cartesianness is visible already on underlying opens: for an
  open \(V \subseteq U_i\),
  \[
    \iota_j^{-1}\bigl(\iota_i(V)\bigr)
      \;=\; (t_{ij} \circ f_{ji})\bigl(f_{ij}^{-1}(V)\bigr),
  \]
  a point of \(U_j\) lands in the image of \(V\) under \(\iota_i\) precisely when it
  comes, via the glue relation, from a point of the overlap lying over \(V\). Hence
  the two composite opens functors
  \[
    (\iota_i)_{!} \circ \iota_j^{-1}
      \;=\;
    f_{ij}^{-1} \circ (t_{ij} \circ f_{ji})_{!}
  \]
  agree on objects, and -- morphisms of the opens poset being proof-irrelevant -- as
  functors. Both \(\iota_i^{*}(\iota_j)_{*}\) and \((f_{ij})_{*}(t_{ij}\circ
  f_{ji})^{*}\) are site-level pushforwards along these two (equal) opens functors, so
  the comparison is assembled from the pushforward-composition isomorphisms together
  with the natural isomorphism transporting one opens functor to the other:
  \[
    \beta_{ij}
      \;=\;
    (\text{pushforwardComp})
      \;\circ\;
    (\text{pushforward of the opens-functor equality})
      \;\circ\;
    (\text{pushforwardCongr})
      \;\circ\;
    (\text{pushforwardComp})^{-1}.
  \]
  The one substantive coherence to check is that the comparison glued from the
  object-level opens equality is natural in the open \(V\): on the section level the
  two restriction maps induced by the equal composite functors coincide because the
  connecting morphisms are the unique arrows of the opens poset between the identified
  opens. This is the pointwise compatibility that completes the \(\mathrm{pushforwardCongr}\)
  factor.
\end{proof}

\section{The restriction morphism and effective descent}
\label{sec:gd_descent}

\begin{lemma}

\leanok
[The chart restriction morphism of the glued sheaf]
  \label{lem:glueRestrictionHom}
  \lean{AlgebraicGeometry.Scheme.Modules.glueRestrictionHom}
  \uses{def:gr_glue_equalizer, def:scheme_modules_glue}
  Let \(\mathcal{M}\) be the glued sheaf of \cref{def:scheme_modules_glue}. The
  \(i\)-th \emph{projection} of \(\mathcal{M}\),
  \[
    \pi_i \;:\; \mathcal{M} \longrightarrow (\iota_i)_{*}\,\mathcal{M}_i,
  \]
  is the equalizer inclusion \(\mathcal{M} \hookrightarrow P\)
  (\cref{def:gr_glue_equalizer}) followed by the \(i\)-th product projection. Its
  adjoint transpose along the chart immersion \(\iota_i\), under the
  pullback--pushforward adjunction \(\iota_i^{*} \dashv (\iota_i)_{*}\), is the
  \emph{restriction morphism}
  \[
    \mathrm{r}_i \;:\; \iota_i^{*}\,\mathcal{M} \longrightarrow \mathcal{M}_i .
  \]
  Equivalently, \(\mathrm{r}_i\) is the unique morphism whose pushforward, precomposed
  with the unit, recovers \(\pi_i\).
\end{lemma}
\begin{proof}
  \uses{def:gr_glue_equalizer, def:scheme_modules_glue}
  \leanok
  The projection \(\pi_i\) is well-defined by \cref{def:gr_glue_equalizer}: the glued
  sheaf is the equalizer \(\operatorname{eq}(a,b) \hookrightarrow P\), and composing
  the inclusion with the \(i\)-th projection of the product \(P = \prod_i
  (\iota_i)_{*}\mathcal{M}_i\) gives a map to \((\iota_i)_{*}\mathcal{M}_i\). The
  restriction morphism is then its image under the natural bijection
  \[
    \operatorname{Hom}_{\mathcal{O}_X}\bigl(\mathcal{M}, (\iota_i)_{*}\mathcal{M}_i\bigr)
      \;\cong\;
    \operatorname{Hom}_{\mathcal{O}_{U_i}}\bigl(\iota_i^{*}\mathcal{M}, \mathcal{M}_i\bigr)
  \]
  supplied by the adjunction \(\iota_i^{*} \dashv (\iota_i)_{*}\). The stated
  equivalent description is the unit--counit form of this transpose.
\end{proof}

\begin{theorem}

\leanok
[Effective descent: the chart restriction morphism is an isomorphism]
  \label{thm:isIso_glueRestrictionHom}
  \lean{AlgebraicGeometry.Scheme.Modules.isIso_glueRestrictionHom}
  \uses{lem:glueRestrictionHom, lem:glueOverlapBaseChangeIso,
    lem:isLimitPullbackCone_mathlib, lem:eq_of_locally_eq_mathlib,
    def:gr_glueRestrictionInv, lem:gr_glueRestrict_hom_ext,
    lem:gr_glueChartComponent_self_counit,
    lem:gr_glueRestrictionHom_glueChartComponent,
    lem:gr_glueRestrictionInv_proj}
  Assume the transition family \((g_{ij})\) satisfies the cocycle conditions (C1) and
  (C2) of \cref{def:scheme_modules_glue}. Then for every chart index \(i\) the
  restriction morphism \(\mathrm{r}_i : \iota_i^{*}\mathcal{M} \to \mathcal{M}_i\)
  (\cref{lem:glueRestrictionHom}) is an isomorphism of sheaves of
  \(\mathcal{O}_{U_i}\)-modules.
\end{theorem}
\begin{proof}
  \uses{lem:glueRestrictionHom, lem:glueOverlapBaseChangeIso,
    lem:isLimitPullbackCone_mathlib, lem:eq_of_locally_eq_mathlib,
    def:gr_glueRestrictionInv, lem:gr_glueRestrict_hom_ext,
    lem:gr_glueChartComponent_self_counit,
    lem:gr_glueRestrictionHom_glueChartComponent,
    lem:gr_glueRestrictionInv_proj}
    \leanok
  The chart pullback \(\iota_i^{*}\) is naturally isomorphic to the site-level
  restriction functor (\cref{lem:modules_restrictFunctorIsoPullback_mathlib}), which
  is a right adjoint; consequently \(\iota_i^{*}\) \emph{preserves limits}. Applying
  it to the descent equalizer \(\mathcal{M} = \operatorname{eq}(a,b)\) exhibits
  \(\iota_i^{*}\mathcal{M}\) as the equalizer of the restricted legs
  \(\iota_i^{*}a, \iota_i^{*}b\), and applying it to the product
  \(P = \prod_j (\iota_j)_{*}\mathcal{M}_j\) exhibits \(\iota_i^{*}P\) as the product
  \(\prod_j \iota_i^{*}(\iota_j)_{*}\mathcal{M}_j\). The overlap base change
  \(\beta_{ij}\) of \cref{lem:glueOverlapBaseChangeIso} then identifies the \(j\)-th
  factor \(\iota_i^{*}(\iota_j)_{*}\mathcal{M}_j\) with
  \((f_{ij})_{*}\bigl((t_{ij} \circ f_{ji})^{*}\mathcal{M}_j\bigr)\). Under these
  identifications the restricted descent data becomes the data attached to the chart
  \(U_i\) by the cover \(\{f_{ij}\}\) of \(U_i\), and the \(i\)-th factor (where
  \(j = i\)) is, by (C1), \(\mathcal{M}_i\) itself.

  We produce an inverse \(\mathrm{s}_i : \mathcal{M}_i \to \iota_i^{*}\mathcal{M}\).
  Because \(\iota_i^{*}\mathcal{M}\) is the equalizer of the restricted legs, it
  suffices to give a morphism \(\mathcal{M}_i \to \iota_i^{*}P\) -- equivalently a
  compatible family of morphisms into the factors -- that equalizes
  \(\iota_i^{*}a\) and \(\iota_i^{*}b\). The \(j\)-th component is
  \[
    \mathcal{M}_i
      \xrightarrow{\ \eta_{f_{ij}}\ }
    (f_{ij})_{*}\,f_{ij}^{*}\mathcal{M}_i
      \xrightarrow{\ (f_{ij})_{*} g_{ij}\ }
    (f_{ij})_{*}\,(t_{ij} \circ f_{ji})^{*}\mathcal{M}_j
      \xrightarrow{\ \beta_{ij}^{-1}\ }
    \iota_i^{*}(\iota_j)_{*}\mathcal{M}_j ,
  \]
  the unit of the pullback--pushforward adjunction along \(f_{ij}\), followed by the
  pushforward of the transition isomorphism \(g_{ij}\), followed by the inverse of the
  overlap base change. That this family equalizes the two restricted legs is exactly
  the triple-overlap multiplicativity (C2), transported through \(\beta\); and the
  composite \(\mathrm{r}_i \circ \mathrm{s}_i = \mathrm{id}_{\mathcal{M}_i}\) reduces,
  via the \(j = i\) component, to the self-identity (C1) together with the counit
  triangle identity of the adjunction. The reverse composite
  \(\mathrm{s}_i \circ \mathrm{r}_i = \mathrm{id}\) holds because both sides are
  equalizer maps with the same components into \(P\); two such agree once their
  factors agree, which is the separatedness of the sheaf equalizer over the chart
  cover. The limit computation underlying these equalities is the sheaf
  Mayer--Vietoris pullback square (\cref{lem:isLimitPullbackCone_mathlib}), and the
  final identification of sections is the uniqueness-over-a-cover principle
  (\cref{lem:eq_of_locally_eq_mathlib}). Hence \(\mathrm{r}_i\) is an isomorphism.
\end{proof}

\section{Construction of the inverse: the descent obligations}
\label{sec:gd_inverse}

This section records, as sublemmas of \cref{thm:isIso_glueRestrictionHom}, the
internal machinery of the proof: the structure-sheaf compatibilities feeding the
overlap base change \(\beta_{ij}\), the bridges between the geometric and the
site-level pullback--pushforward adjunctions, the explicit construction of the
candidate inverse \(\mathrm{s}_i\) as an equalizer lift, and the three obligations
that consume the cocycle conditions (C1)/(C2). Every block carries a one-to-one Lean
declaration.

\subsection{Section-level compatibilities of the overlap square}

\begin{lemma}\leanok
[Transport of \(\mathrm{appLE}\) along equal morphisms]
  \label{lem:gr_appLE_congr_mor}
  \lean{AlgebraicGeometry.Scheme.Modules.appLE_congr_mor}
  Let \(f, g : A \to B\) be morphisms of schemes with \(f = g\), let \(U\) be an open of
  \(B\) and \(W\) an open of \(A\) with \(W \subseteq f^{-1}(U)\) (equivalently
  \(W \subseteq g^{-1}(U)\)). Then the two induced section maps
  \(\Gamma(B, U) \to \Gamma(A, W)\) along \(f\) and along \(g\) coincide:
  \(f^{\mathrm{LE}}_{U,W} = g^{\mathrm{LE}}_{U,W}\).
\end{lemma}
\begin{proof}

\leanok
  Substituting the equality \(f = g\) reduces both sides to the same map; the two
  open-inclusion witnesses \(W \subseteq f^{-1}(U)\) and \(W \subseteq g^{-1}(U)\) are
  proof-irrelevant.
\end{proof}

\begin{lemma}\leanok
[Structure-sheaf compatibility of the chart--overlap square]
  \label{lem:gr_overlap_appIso_compat}
  \lean{AlgebraicGeometry.Scheme.Modules.glueData_overlap_appIso_compat}
  \uses{lem:gr_appLE_congr_mor, def:scheme_modules_glue}
  For chart indices \(i, j\) and an open \(V \subseteq U_i\), the two composite
  section maps
  \[
    \Gamma(U_i, V)\;\longrightarrow\;
      \Gamma\bigl(U_j,\ (t_{ij}\circ f_{ji})\bigl(f_{ij}^{-1}(V)\bigr)\bigr)
  \]
  of the chart--overlap square agree. The first goes ``through the glued scheme'':
  the inverse of the comparison isomorphism of \(\iota_i\) on \(V\), the section map of
  \(\iota_j\) over \(\iota_i(V)\), and the restriction along the opens identity
  \(\iota_j^{-1}(\iota_i(V)) = (t_{ij}\circ f_{ji})(f_{ij}^{-1}(V))\). The second goes
  ``through the overlap'': the section map of \(f_{ij}\) over \(V\) followed by the
  inverse comparison isomorphism of \(t_{ij}\circ f_{ji}\).
\end{lemma}
\begin{proof}
  \uses{lem:gr_appLE_congr_mor, def:scheme_modules_glue}
  \leanok
  Both composites are \(\mathrm{appLE}\)s of the two composite morphisms of the square,
  which are equal by the glue condition \((t_{ij}\circ f_{ji})\circ \iota_j =
  f_{ij}\circ \iota_i\) of \cref{def:scheme_modules_glue}. Moving the two comparison
  isomorphisms to \(\mathrm{appLE}\) form (using \(\mathrm{appIso}\) as an
  \(\mathrm{appLE}\), the composition law \(g^{\mathrm{LE}}\circ f^{\mathrm{LE}} =
  (f\circ g)^{\mathrm{LE}}\), and the absorption of presheaf restrictions into the
  open-inclusion witness) reduces the claim to the equality of the two
  \(\mathrm{appLE}\)s of the equal composites, which is \cref{lem:gr_appLE_congr_mor}.
\end{proof}

\begin{lemma}\leanok
[Sections of the overlap base change, inverse side]
  \label{lem:gr_overlapBaseChange_inv_app}
  \lean{AlgebraicGeometry.Scheme.Modules.glueOverlapBaseChangeIso_inv_app_app}
  \uses{lem:glueOverlapBaseChangeIso}
  For a sheaf \(\mathcal{N}\) on \(U_j\) and an open \(V \subseteq U_i\), the inverse
  of the overlap base change \(\beta_{ij}^{-1}\), evaluated on sections over \(V\), is
  the restriction map of \(\mathcal{N}\) along the opens identity
  \(\iota_j^{-1}(\iota_i(V)) = (t_{ij}\circ f_{ji})(f_{ij}^{-1}(V))\) of
  \cref{lem:glueOverlapBaseChangeIso}:
  \[
    (\beta_{ij}^{-1})_{\mathcal{N}}\big|_V
      \;=\;
    \mathcal{N}\!\left(\,(t_{ij}\circ f_{ji})(f_{ij}^{-1}V) \;=\; \iota_j^{-1}(\iota_i V)\,\right).
  \]
\end{lemma}
\begin{proof}
  \uses{lem:glueOverlapBaseChangeIso}
  \leanok
  Each of the four functorial factors of \(\beta_{ij}\) (\cref{lem:glueOverlapBaseChangeIso})
  acts on sections over \(V\) either as the identity or as a single presheaf
  restriction along the opens identity; composing them leaves exactly the named
  restriction map. The verification is pointwise on sections.
\end{proof}

\begin{lemma}\leanok
[Sections of the overlap base change, forward side]
  \label{lem:gr_overlapBaseChange_hom_app}
  \lean{AlgebraicGeometry.Scheme.Modules.glueOverlapBaseChangeIso_hom_app_app}
  \uses{lem:glueOverlapBaseChangeIso}
  Symmetrically, the forward overlap base change \(\beta_{ij}\), evaluated on sections
  over an open \(V \subseteq U_i\), is the restriction map of \(\mathcal{N}\) along the
  reverse opens identity \(\iota_j^{-1}(\iota_i(V)) = (t_{ij}\circ f_{ji})(f_{ij}^{-1}(V))\).
\end{lemma}
\begin{proof}
  \uses{lem:glueOverlapBaseChangeIso}
  \leanok
  Identical to \cref{lem:gr_overlapBaseChange_inv_app} with the roles of source and
  target opens exchanged: pointwise on sections, the composite of the four factors is
  the single presheaf restriction along the (reversed) opens identity.
\end{proof}

\subsection{Bridges between the geometric and the site-level adjunction}

The chart pullback \(\iota^{*}\) is built by sheafification and carries the
\emph{geometric} adjunction \(\iota^{*} \dashv \iota_{*}\); along an open immersion the
\emph{site-level} restriction functor \(\iota^{\mathrm{res}}\) carries a second
adjunction with concrete unit and counit. They share the right adjoint \(\iota_{*}\),
and \cref{lem:modules_restrictFunctorIsoPullback_mathlib} is the uniqueness comparison
of left adjoints between them. The next blocks transport units, counits and
congruence casts across that comparison, so that the descent obligations can be
discharged by section-level computations.

\begin{lemma}\leanok
[Unit comparison across the adjunctions]
  \label{lem:gr_restrictAdjunction_unit_iso}
  \lean{AlgebraicGeometry.Scheme.Modules.restrictAdjunction_unit_app_iso}
  \uses{lem:modules_restrictFunctorIsoPullback_mathlib}
  For an open immersion \(f\) and a sheaf \(\mathcal{N}\) on the target, the geometric
  unit is the site-level unit followed by the pushforward of the left-adjoint
  comparison isomorphism:
  \[
    \eta^{\mathrm{res}}_{\mathcal{N}}\ \circ\
      f_{*}\bigl((\text{comparison})_{\mathcal{N}}\bigr)
    \;=\; \eta^{\mathrm{geo}}_{\mathcal{N}}.
  \]
\end{lemma}
\begin{proof}
  \uses{lem:modules_restrictFunctorIsoPullback_mathlib}
  \leanok
  This is the general identity relating the unit of an adjunction to the unit of any
  other adjunction with the same right adjoint, transported by the uniqueness
  isomorphism of left adjoints (\cref{lem:modules_restrictFunctorIsoPullback_mathlib}).
\end{proof}

\begin{lemma}\leanok
[Counit comparison across the adjunctions]
  \label{lem:gr_restrictIsoPullback_counit}
  \lean{AlgebraicGeometry.Scheme.Modules.restrictFunctorIsoPullback_hom_app_counit}
  \uses{lem:modules_restrictFunctorIsoPullback_mathlib}
  Dually, the left-adjoint comparison isomorphism on \(f_{*}\mathcal{N}\) followed by
  the geometric counit equals the site-level counit:
  \[
    (\text{comparison})_{f_{*}\mathcal{N}}\ \circ\ \varepsilon^{\mathrm{geo}}_{\mathcal{N}}
      \;=\; \varepsilon^{\mathrm{res}}_{\mathcal{N}}.
  \]
\end{lemma}
\begin{proof}
  \uses{lem:modules_restrictFunctorIsoPullback_mathlib}
  \leanok
  The dual of \cref{lem:gr_restrictAdjunction_unit_iso}: the counit of one adjunction
  is conjugated to the counit of the other by the uniqueness comparison of left
  adjoints.
\end{proof}

\begin{lemma}\leanok
[Congruence cast at the pullback level]
  \label{lem:gr_pullbackCongr_hom_eqToHom}
  \lean{AlgebraicGeometry.Scheme.Modules.pullbackCongr_hom_app_eqToHom}
  For equal morphisms \(\varphi = \psi : X \to Y\) and a sheaf \(\mathcal{N}\) on \(Y\),
  the congruence isomorphism \(\varphi^{*}\mathcal{N} \cong \psi^{*}\mathcal{N}\)
  induced by \(\varphi = \psi\) is the canonical transport (the \(\mathrm{eqToHom}\) of
  the object-level equality \(\varphi^{*}\mathcal{N} = \psi^{*}\mathcal{N}\)).
\end{lemma}
\begin{proof}

\leanok
  Substituting \(\varphi = \psi\) reduces the congruence cast to the identity, which is
  the canonical transport of a reflexivity.
\end{proof}

\begin{lemma}\leanok
[Trivial site-level congruence cast]
  \label{lem:gr_restrictFunctorCongr_rfl}
  \lean{AlgebraicGeometry.Scheme.Modules.restrictFunctorCongr_rfl_hom_app}
  For an open immersion \(\varphi\) the site-level congruence cast associated to the
  reflexivity \(\varphi = \varphi\) is the identity natural transformation.
\end{lemma}
\begin{proof}

\leanok
  Pointwise on sections the cast is the presheaf restriction along a self-equality of
  opens, which is the identity by functoriality.
\end{proof}

\begin{lemma}\leanok
[Comparison intertwines the two congruence casts]
  \label{lem:gr_restrictIsoPullback_congr}
  \lean{AlgebraicGeometry.Scheme.Modules.restrictFunctorIsoPullback_congr}
  \uses{lem:gr_pullbackCongr_hom_eqToHom, lem:gr_restrictFunctorCongr_rfl,
    lem:modules_restrictFunctorIsoPullback_mathlib}
  For equal open immersions \(\varphi = \psi : X \to Y\) the left-adjoint comparison
  isomorphism intertwines the site-level and the pullback-level congruence casts:
  \[
    (\text{comparison}_\varphi)_{\mathcal{N}}\ \circ\
      (\varphi^{*}\!=\psi^{*})_{\mathcal{N}}
    \;=\;
    (\varphi^{\mathrm{res}}\!=\psi^{\mathrm{res}})_{\mathcal{N}}\ \circ\
      (\text{comparison}_\psi)_{\mathcal{N}}.
  \]
\end{lemma}
\begin{proof}
  \uses{lem:gr_pullbackCongr_hom_eqToHom, lem:gr_restrictFunctorCongr_rfl}
  \leanok
  Substituting \(\varphi = \psi\) collapses both congruence casts to identities
  (\cref{lem:gr_pullbackCongr_hom_eqToHom}, \cref{lem:gr_restrictFunctorCongr_rfl}),
  and the two comparison isomorphisms then coincide.
\end{proof}

\subsection{The candidate inverse}

\begin{definition}\leanok
[Overlap base change in pullback form]
  \label{def:gr_glueOverlapFactorIso}
  \lean{AlgebraicGeometry.Scheme.Modules.glueOverlapFactorIso}
  \uses{lem:glueOverlapBaseChangeIso, lem:modules_restrictFunctorIsoPullback_mathlib}
  Conjugating the overlap base change \(\beta_{ij}\)
  (\cref{lem:glueOverlapBaseChangeIso}) on both sides by the comparison
  \cref{lem:modules_restrictFunctorIsoPullback_mathlib} between the site-level
  restriction and the geometric pullback yields the isomorphism, in geometric form,
  \[
    \gamma_{ij} \;:\;
      \iota_i^{*}\bigl((\iota_j)_{*}\,\mathcal{M}_j\bigr)
        \;\xrightarrow{\ \sim\ }\;
      (f_{ij})_{*}\bigl((t_{ij}\circ f_{ji})^{*}\,\mathcal{M}_j\bigr),
  \]
  evaluated at the input sheaf \(\mathcal{M}_j\).
\end{definition}

\begin{definition}\leanok
[The \(j\)-th component of the candidate inverse]
  \label{def:gr_glueChartComponent}
  \lean{AlgebraicGeometry.Scheme.Modules.glueChartComponent}
  \uses{def:gr_glueOverlapFactorIso}
  Given the transition isomorphisms \(g_{ij} : f_{ij}^{*}\mathcal{M}_i \xrightarrow{\sim}
  (t_{ij}\circ f_{ji})^{*}\mathcal{M}_j\), the \(j\)-th component of the candidate
  inverse \(\mathrm{s}_i\) is the composite
  \[
    \mathcal{M}_i
      \xrightarrow{\ \eta_{f_{ij}}\ }
    (f_{ij})_{*}\,f_{ij}^{*}\mathcal{M}_i
      \xrightarrow{\ (f_{ij})_{*}g_{ij}\ }
    (f_{ij})_{*}\,(t_{ij}\circ f_{ji})^{*}\mathcal{M}_j
      \xrightarrow{\ \gamma_{ij}^{-1}\ }
    \iota_i^{*}\bigl((\iota_j)_{*}\mathcal{M}_j\bigr),
  \]
  i.e. the geometric unit along \(f_{ij}\), the pushforward of the transition
  \(g_{ij}\), and the inverse of \(\gamma_{ij}\) (\cref{def:gr_glueOverlapFactorIso}).
\end{definition}

\begin{definition}\leanok
[The candidate-inverse family]
  \label{def:gr_glueChartFamily}
  \lean{AlgebraicGeometry.Scheme.Modules.glueChartFamily}
  \uses{def:gr_glueChartComponent, def:gr_glue_equalizer}
  Assembling the components \cref{def:gr_glueChartComponent} through the
  product-preservation comparison of \(\iota_i^{*}\) on the pushforward product
  \(P = \prod_j (\iota_j)_{*}\mathcal{M}_j\) (\cref{def:gr_glue_equalizer}) gives a
  single morphism
  \[
    \mathcal{M}_i \;\longrightarrow\; \iota_i^{*}P,
  \]
  the candidate inverse before it is corrected to land in the equalizer
  \(\iota_i^{*}\mathcal{M}\).
\end{definition}

\subsection{The triple-overlap toolkit}

The first descent obligation (\cref{lem:gr_glueChartFamily_equalizes}) reduces, factor
by factor of the restricted overlap product, to a single equation over a \emph{triple
overlap}. This subsection records that machinery. Fix three chart indices \(i, p, q\)
and form the triple overlap
\[
  V_{ipq} \;=\; V_{ip} \times_{U_i} V_{iq},
\]
with first projection \(\mathrm{fst} : V_{ipq} \to V_{ip}\). Two morphisms out of
\(V_{ipq}\) organise everything: the \emph{chart leg}
\(q_{pq} = \mathrm{fst}\circ f_{ip} : V_{ipq} \to U_i\) (down to \(U_i\)), and the
\emph{transport leg} \(\tau = t'_{ipq}\circ \mathrm{fst}_{pq} : V_{ipq} \to V_{pq}\)
(across to the pair overlap, through the cocycle transport \(t'_{ipq}\)). The square
\[
  \begin{array}{ccc}
    V_{ipq} & \xrightarrow{\ \ \tau\ \ } & V_{pq} \\
    \ \downarrow{\scriptstyle q_{pq}} & & \downarrow{\scriptstyle f_{pq}\circ \iota_p} \\
    U_i & \xrightarrow{\ \ \iota_i\ \ } & X
  \end{array}
\]
is the triple analogue of the chart--overlap square of \cref{lem:glueOverlapBaseChangeIso}.

\begin{lemma}

\leanok
[Triple-overlap bridges of the glue datum]
  \label{lem:gr_glueData_bridges}
  \lean{AlgebraicGeometry.Scheme.Modules.glueData_bridge_src,
    AlgebraicGeometry.Scheme.Modules.glueData_bridge_mid,
    AlgebraicGeometry.Scheme.Modules.glueData_bridge_tgt}
  \uses{def:scheme_modules_glue}
  On the triple overlap \(V_{ijk} = V_{ij}\times_{U_i} V_{ik}\) the following three
  identities of morphisms hold, lining up the endpoints of the three base-change
  transports of the cocycle equation (C2):
  \begin{enumerate}
    \item \emph{(source)} the two projections of \(V_{ijk}\) to \(V_{ij}\) and \(V_{ik}\),
      followed by the overlap immersions \(f_{ij}\) and \(f_{ik}\) into \(U_i\),
      agree;
    \item \emph{(middle)} the \(ij\)-transport's target leg
      \((t_{ij}\circ f_{ji})\) and the \(jk\)-transport's source leg
      (through the cocycle transport \(t'_{ijk}\) and \(f_{jk}\)) agree as morphisms
      to \(U_j\);
    \item \emph{(target)} the composite \(jk\)-after-\(ij\) transport's target leg
      \((t_{jk}\circ f_{kj})\) and the \(ik\)-transport's target leg
      \((t_{ik}\circ f_{ki})\) agree as morphisms to \(U_k\).
  \end{enumerate}
\end{lemma}
\begin{proof}
  \leanok
  \uses{def:scheme_modules_glue}
  The source identity is the universal property (pullback condition) of \(V_{ijk}\).
  The middle identity follows from the transition factorisation
  \(t'_{ijk}\) of \cref{def:scheme_modules_glue} together with the pullback condition.
  The target identity -- the heart of the cocycle -- combines the same transition
  factorisation, the pullback condition, the involutivity \(t_{ik}\circ t_{ki}=\mathrm{id}\)
  and the cocycle condition of \cref{def:scheme_modules_glue}.
\end{proof}

\begin{lemma}

\leanok
[Commutativity of the triple-overlap square]
  \label{lem:gr_glueData_triple_square}
  \lean{AlgebraicGeometry.Scheme.Modules.glueData_triple_square}
  \uses{lem:gr_glueData_bridges, def:scheme_modules_glue}
  The triple-overlap square above commutes: the transport leg \(\tau\) followed by
  \(f_{pq}\circ \iota_p\) equals the chart leg \(q_{pq}\) followed by \(\iota_i\), as
  morphisms \(V_{ipq}\to X\).
\end{lemma}
\begin{proof}
  \uses{lem:gr_glueData_bridges, def:scheme_modules_glue}
  \leanok
  Substitute the middle bridge \cref{lem:gr_glueData_bridges} to rewrite the transport
  leg, then apply the glue condition \((t_{ip}\circ f_{pi})\circ\iota_p =
  f_{ip}\circ\iota_i\) of \cref{def:scheme_modules_glue}.
\end{proof}

\begin{lemma}

\leanok
[Triple-overlap opens identity]
  \label{lem:gr_glueData_triple_preimage}
  \lean{AlgebraicGeometry.Scheme.Modules.glueData_preimage_image_eq₃}
  \uses{lem:gr_glueData_bridges, lem:gr_glueData_triple_square, def:scheme_modules_glue}
  For an open \(V \subseteq U_i\), the preimage under \(f_{pq}\circ\iota_p\) of the image
  \(\iota_i(V)\) in the glued scheme equals the image, under the transport leg \(\tau\),
  of the preimage \(q_{pq}^{-1}(V)\):
  \[
    (f_{pq}\circ\iota_p)^{-1}\bigl(\iota_i(V)\bigr)
      \;=\; \tau\bigl(q_{pq}^{-1}(V)\bigr).
  \]
  This is the underlying-opens cartesianness of the triple-overlap square.
\end{lemma}
\begin{proof}
  \uses{lem:gr_glueData_bridges, lem:gr_glueData_triple_square, def:scheme_modules_glue}
  \leanok
  A point of \(V_{pq}\) lands in \(\iota_i(V)\) precisely when it is glued, through the
  overlap \(V_{ip}\), to a point of \(U_i\) lying in \(V\); the gluing-incidence
  criterion at the pairs \((i,p)\) and \((i,q)\) certifies that such a point lifts to
  \(V_{ipq}\), and injectivity of the open immersion \(f_{pq}\) pins down the transport,
  using the middle bridge \cref{lem:gr_glueData_bridges}. The reverse inclusion is the
  commutativity \cref{lem:gr_glueData_triple_square}.
\end{proof}

\begin{lemma}

\leanok
[The two composite opens functors of the triple square agree]
  \label{lem:gr_glueData_triple_opensFunctor_eq}
  \lean{AlgebraicGeometry.Scheme.Modules.glueData_triple_opensFunctor_eq}
  \uses{lem:gr_glueData_triple_preimage}
  The composite opens functor ``into the glued scheme along \(\iota_i\), then back to
  \(V_{pq}\) along \(f_{pq}\circ\iota_p\)'' equals the composite ``down to \(V_{ipq}\)
  along \(q_{pq}\), then across to \(V_{pq}\) along \(\tau\)''.
\end{lemma}
\begin{proof}
  \uses{lem:gr_glueData_triple_preimage}
  \leanok
  The two functors agree on objects by the opens identity
  \cref{lem:gr_glueData_triple_preimage}, and on morphisms because the connecting arrows
  of the opens poset are proof-irrelevant.
\end{proof}

\begin{lemma}

\leanok
[Structure-sheaf compatibility of the triple square]
  \label{lem:gr_glueData_triple_appIso_compat}
  \lean{AlgebraicGeometry.Scheme.Modules.glueData_triple_appIso_compat}
  \uses{lem:gr_appLE_congr_mor, lem:gr_glueData_triple_square}
  For an open \(V \subseteq U_i\) the two composite section maps
  \(\Gamma(U_i, V) \to \Gamma\bigl(V_{pq}, \tau(q_{pq}^{-1}(V))\bigr)\) of the triple
  square -- ``through the glued scheme'' and ``through the triple overlap'' -- coincide.
\end{lemma}
\begin{proof}
  \uses{lem:gr_appLE_congr_mor, lem:gr_glueData_triple_square}
  \leanok
  Both composites are the \(\mathrm{appLE}\) of the two (equal, by
  \cref{lem:gr_glueData_triple_square}) composite morphisms of the square; the equality
  is then \cref{lem:gr_appLE_congr_mor}, exactly as in the pair-level
  \cref{lem:gr_overlap_appIso_compat}.
\end{proof}

\begin{lemma}

\leanok
[Triple-overlap base change (\(\beta_{ipq}\))]
  \label{lem:gr_glueTripleBaseChangeIso}
  \lean{AlgebraicGeometry.Scheme.Modules.glueTripleBaseChangeIso}
  \uses{lem:gr_glueData_triple_preimage, lem:gr_glueData_triple_opensFunctor_eq,
    lem:gr_glueData_triple_appIso_compat}
  There is a natural isomorphism of functors
  \(\ShMod(\mathcal{O}_{V_{pq}}) \to \ShMod(\mathcal{O}_{V_{ipq}})\),
  \[
    \beta_{ipq} \;:\;
      (f_{pq}\circ\iota_p)_{*}\circ \iota_i^{*}
        \;\xrightarrow{\ \sim\ }\;
      \tau^{*}\circ (q_{pq})_{*},
  \]
  identifying the chart restriction of a pushforward from the pair overlap \(V_{pq}\)
  with the pushforward along \(q_{pq}\) of the restriction along \(\tau\). It is
  assembled by the same four-factor recipe as the pair-level overlap base change
  (\cref{lem:glueOverlapBaseChangeIso}) over the opens identity
  \cref{lem:gr_glueData_triple_preimage}.
\end{lemma}
\begin{proof}
  \uses{lem:gr_glueData_triple_preimage, lem:gr_glueData_triple_opensFunctor_eq,
    lem:gr_glueData_triple_appIso_compat}
    \leanok
  Both sides are site-level pushforwards along the two equal composite opens functors
  \cref{lem:gr_glueData_triple_opensFunctor_eq}; the comparison is the pushforward-com\-position
  isomorphisms together with the transport of one opens functor to the other, with the
  section-level coherence supplied by \cref{lem:gr_glueData_triple_appIso_compat}.
\end{proof}

\begin{lemma}

\leanok
[Sections of the triple base change, inverse side]
  \label{lem:gr_glueTripleBaseChangeIso_inv_app}
  \lean{AlgebraicGeometry.Scheme.Modules.glueTripleBaseChangeIso_inv_app_app}
  \uses{lem:gr_glueTripleBaseChangeIso, lem:gr_glueData_triple_preimage}
  On sections over an open \(V\subseteq U_i\), the inverse \(\beta_{ipq}^{-1}\) is the
  restriction map of the sheaf along the opens identity
  \cref{lem:gr_glueData_triple_preimage}.
\end{lemma}
\begin{proof}
  \uses{lem:gr_glueTripleBaseChangeIso, lem:gr_glueData_triple_preimage}
  \leanok
  Each functorial factor of \(\beta_{ipq}\) acts on sections as the identity or as a
  single presheaf restriction along the opens identity; composing leaves exactly the
  named restriction. Pointwise on sections.
\end{proof}

\begin{lemma}

\leanok
[Sections of the triple base change, forward side]
  \label{lem:gr_glueTripleBaseChangeIso_hom_app}
  \lean{AlgebraicGeometry.Scheme.Modules.glueTripleBaseChangeIso_hom_app_app}
  \uses{lem:gr_glueTripleBaseChangeIso, lem:gr_glueData_triple_preimage}
  Symmetrically, on sections over \(V\subseteq U_i\) the forward \(\beta_{ipq}\) is the
  restriction map along the reverse opens identity
  \cref{lem:gr_glueData_triple_preimage}.
\end{lemma}
\begin{proof}
  \uses{lem:gr_glueTripleBaseChangeIso, lem:gr_glueData_triple_preimage}
  \leanok
  Identical to \cref{lem:gr_glueTripleBaseChangeIso_inv_app} with source and target
  opens exchanged.
\end{proof}

\begin{lemma}

\leanok
[Triple-overlap base change in pullback form]
  \label{lem:gr_glueTripleFactorIso}
  \lean{AlgebraicGeometry.Scheme.Modules.glueTripleFactorIso}
  \uses{lem:gr_glueTripleBaseChangeIso, lem:modules_restrictFunctorIsoPullback_mathlib}
  For a sheaf \(\mathcal{N}\) on the pair overlap \(V_{pq}\), conjugating \(\beta_{ipq}\)
  on both sides by the comparison between site-level restriction and geometric pullback
  (\cref{lem:modules_restrictFunctorIsoPullback_mathlib}) yields the geometric-form
  isomorphism
  \[
    \iota_i^{*}\bigl((f_{pq}\circ\iota_p)_{*}\mathcal{N}\bigr)
      \;\xrightarrow{\ \sim\ }\;
    (q_{pq})_{*}\bigl(\tau^{*}\mathcal{N}\bigr).
  \]
\end{lemma}
\begin{proof}
  \uses{lem:gr_glueTripleBaseChangeIso, lem:modules_restrictFunctorIsoPullback_mathlib}
  \leanok
  Conjugate \(\beta_{ipq}\) (\cref{lem:gr_glueTripleBaseChangeIso}) by the left-adjoint
  comparison \cref{lem:modules_restrictFunctorIsoPullback_mathlib} on the chart factor
  and by the same comparison pushed forward along \(q_{pq}\) on the transport factor.
\end{proof}

\begin{lemma}

\leanok
[Transpose of the triple base change (mate core)]
  \label{lem:gr_glueTripleFactor_transpose}
  \lean{AlgebraicGeometry.Scheme.Modules.glueTripleFactor_transpose}
  \uses{lem:gr_glueTripleFactorIso, lem:gr_glueTripleBaseChangeIso,
    lem:gr_restrictAdjunction_unit_iso, lem:gr_glueData_triple_preimage,
    lem:gr_glueData_triple_square}
  The adjoint transpose along \(\iota_i\) of the pullback-form triple base change
  \cref{lem:gr_glueTripleFactorIso} is the canonical four-functor comparison of the
  triple square: the unit along \(\tau\) pushed forward along \(f_{pq}\circ\iota_p\),
  regrouped by the pushforward-composition comparison, reindexed by the
  pushforward-congruence along \cref{lem:gr_glueData_triple_square}, and ungrouped by the
  inverse pushforward-composition. The identity is free of the cocycle hypotheses.
\end{lemma}
\begin{proof}
  \uses{lem:gr_glueTripleFactorIso, lem:gr_glueTripleBaseChangeIso,
    lem:gr_restrictAdjunction_unit_iso, lem:gr_glueData_triple_preimage}
    \leanok
  Verbatim the proof of the pair-level \cref{lem:gr_glueOverlapFactor_transpose} with the
  substitutions \(\iota_j \rightsquigarrow f_{pq}\circ\iota_p\),
  \(t_{ij}\circ f_{ji}\rightsquigarrow \tau\), \(f_{ij}\rightsquigarrow q_{pq}\): the
  comparison isomorphisms conjugating the chart factor cancel against unit/counit by
  \cref{lem:gr_restrictAdjunction_unit_iso}, leaving \(\beta_{ipq}\), whose section form
  (\cref{lem:gr_glueData_triple_preimage}) folds the residual cycle of restrictions to
  the stated four-factor composite.
\end{proof}

\begin{lemma}

\leanok
[Mate recognition of the triple factor transpose]
  \label{lem:gr_glueTripleFactor_mate}
  \lean{AlgebraicGeometry.Scheme.Modules.glueTripleFactor_mate}
  \uses{lem:gr_glueTripleFactor_transpose, lem:gr_homEquiv_conjugateEquiv_app,
    lem:conjugateEquiv_pullbackComp_inv_mathlib, lem:gr_glueData_triple_square}
  For any morphism \(m : \mathcal{W} \to (f_{pq}\circ\iota_p)_{*}\mathcal{N}\) on the
  glued scheme, the conjugated transport of the inverse transpose of \(m\) along the
  triple square equals the inverse transpose, along \(q_{pq}\), of \(\iota_i^{*}(m)\)
  followed by the pullback-form triple base change \cref{lem:gr_glueTripleFactorIso}. In
  other words, \cref{lem:gr_glueTripleFactor_transpose} is the mate of
  \cref{lem:gr_glueTripleFactorIso} through the pullback pseudofunctor.
\end{lemma}
\begin{proof}
  \uses{lem:gr_glueTripleFactor_transpose, lem:gr_homEquiv_conjugateEquiv_app,
    lem:conjugateEquiv_pullbackComp_inv_mathlib}
    \leanok
  By injectivity of the two transpose bijections compare the double transposes;
  naturality of the transpose in its left argument peels \(m\) off, the conjugate of the
  pullback-composition isomorphism along the triple square
  (\cref{lem:gr_homEquiv_conjugateEquiv_app},
  \cref{lem:conjugateEquiv_pullbackComp_inv_mathlib}) is the inverse
  pushforward-composition isomorphism, and folding the unit pair reduces the claim to the
  mate core \cref{lem:gr_glueTripleFactor_transpose} -- the same proof as the pair-level
  \cref{lem:gr_glueOverlapFactor_mate}.
\end{proof}

\begin{lemma}

\leanok
[Transpose of the first descent-leg factor]
  \label{lem:gr_glueLegA_component_transpose}
  \lean{AlgebraicGeometry.Scheme.Modules.glueLegA_component_transpose}
  \uses{def:gr_glue_equalizer}
  Transposed along the composite overlap immersion \(f_{pq}\circ\iota_p\), the first
  descent-leg factor of \cref{def:gr_glue_equalizer} (the unit along \(f_{pq}\) followed
  by the pushforward-composition comparison) becomes the pullback-pseudofunctor cast
  followed by the chart pullback of the geometric counit.
\end{lemma}
\begin{proof}
  \uses{def:gr_glue_equalizer}
  \leanok
  The transpose of the unit pair at \((f_{pq}, \iota_p)\) is computed by the
  composite-unit identity; the transpose of the counit is the identity by the right
  triangle identity of the adjunction, which consumes the inner factor.
\end{proof}

\begin{lemma}

\leanok
[Transpose of the second descent-leg factor]
  \label{lem:gr_glueLegB_component_transpose}
  \lean{AlgebraicGeometry.Scheme.Modules.glueLegB_component_transpose}
  \uses{def:gr_glue_equalizer, def:scheme_modules_glue}
  Transposed along \(f_{pq}\circ\iota_p\), the second descent-leg factor of
  \cref{def:gr_glue_equalizer} (the unit along \(t_{pq}\circ f_{qp}\), the
  pushforward-composition comparison, the pushforward of \(g_{pq}^{-1}\), and the
  pushforward-congruence along the glue condition) becomes the pullback-pseudofunctor
  casts, the chart pullback of the geometric counit, and the transition \(g_{pq}^{-1}\).
\end{lemma}
\begin{proof}
  \uses{def:gr_glue_equalizer, def:scheme_modules_glue}
  \leanok
  As in \cref{lem:gr_glueLegA_component_transpose}, with two extra moves: naturality of
  the transpose in its right argument absorbs \(g_{pq}^{-1}\), and the composite-congruence
  identity re-indexes along the glue condition of \cref{def:scheme_modules_glue}.
\end{proof}

\begin{lemma}

\leanok
[The candidate-inverse family recovers its components]
  \label{lem:gr_glueChartFamily_pullback_map_pi}
  \lean{AlgebraicGeometry.Scheme.Modules.glueChartFamily_pullback_map_π}
  \uses{def:gr_glueChartFamily, def:gr_glueChartComponent}
  The candidate-inverse family \cref{def:gr_glueChartFamily} followed by the chart
  pullback of the \(j\)-th product projection recovers the \(j\)-th component
  \cref{def:gr_glueChartComponent}:
  \[
    \mathrm{s}_i\ \circ\ \iota_i^{*}(\mathrm{pr}_j) \;=\; \mathrm{s}_i^{(j)}.
  \]
\end{lemma}
\begin{proof}
  \uses{def:gr_glueChartFamily, def:gr_glueChartComponent}
  \leanok
  Unfold the family \cref{def:gr_glueChartFamily} as a product lift composed with the
  inverse product-preservation comparison; the comparison cancels against the lift,
  leaving the defining \(j\)-th component \cref{def:gr_glueChartComponent}.
\end{proof}

\begin{lemma}

\leanok
[Pair component of the equalizing condition (C2 over the triple overlap)]
  \label{lem:gr_glueChartComponent_leg_compat}
  \lean{AlgebraicGeometry.Scheme.Modules.glueChartComponent_leg_compat}
  \uses{def:gr_glueChartComponent, def:gr_glueOverlapFactorIso,
    lem:gr_glueTripleFactorIso, lem:gr_glueTripleFactor_mate,
    lem:gr_glueLegA_component_transpose, lem:gr_glueLegB_component_transpose,
    lem:gr_glueOverlapFactor_mate, lem:gr_glueData_bridges,
    lem:modules_pullback_basechange_transport, lem:gr_pullbackCongr_hom_eqToHom,
    def:scheme_modules_glue}
  For chart indices \(i, p, q\), the \(p\)-th candidate-inverse component
  \cref{def:gr_glueChartComponent} followed by the chart pullback of the first
  descent-leg factor equals the \(q\)-th component followed by the chart pullback of the
  second descent-leg factor. This is the \((p,q)\)-component of the first descent
  obligation, where the triple-overlap multiplicativity (C2) of
  \cref{def:scheme_modules_glue} at \((i,p,q)\) is consumed.
\end{lemma}
\begin{proof}
  \uses{def:gr_glueChartComponent, def:gr_glueOverlapFactorIso,
    lem:gr_glueTripleFactorIso, lem:gr_glueTripleFactor_mate,
    lem:gr_glueLegA_component_transpose, lem:gr_glueLegB_component_transpose,
    lem:gr_glueOverlapFactor_mate, lem:gr_glueData_bridges,
    lem:modules_pullback_basechange_transport, lem:gr_pullbackCongr_hom_eqToHom,
    def:scheme_modules_glue}
    \leanok
  \emph{Reduction.} Cancel the triple base change in pullback form
  \cref{lem:gr_glueTripleFactorIso} (an isomorphism) from the right and transpose the
  equation along the chart leg \(q_{pq}\) (the transpose is a bijection). On both sides
  fire the triple mate \cref{lem:gr_glueTripleFactor_mate} and substitute the leg
  transposes \cref{lem:gr_glueLegA_component_transpose},
  \cref{lem:gr_glueLegB_component_transpose}. This reduces the claim to the
  \emph{conjugated cocycle equation} over \(V_{ipq}\): the chart pullbacks of the two
  components, base-change--transported and re-indexed, agree.

  \emph{The conjugated cocycle.} Expand each component
  \[
    \mathrm{s}_i^{(p)}
      \;=\; \eta_{f_{ip}}\ \circ\ (f_{ip})_{*}g_{ip}\ \circ\ \gamma_{ip}^{-1}
  \]
  (\cref{def:gr_glueChartComponent}) and pull it back along \(q_{pq} = \mathrm{fst}\circ
  f_{ip}\). Three moves clear the remaining factors:
  \begin{enumerate}
    \item the chart pullback \(q_{pq}^{*}\) of the inverse overlap factor
      \(\gamma_{ip}^{-1}\) (\cref{def:gr_glueOverlapFactorIso}) is converted into
      pullback-composition casts and transposes by the \emph{pair} mate
      \cref{lem:gr_glueOverlapFactor_mate} at \((i,p)\);
    \item the chart pullback of the unit \(\eta_{f_{ip}}\) along the morphism
      \(q_{pq}\), which factors through \(f_{ip}\), is consumed by the triangle identity
      of the adjunction threaded through the pullback-composition cast at
      \((\mathrm{fst}, f_{ip})\) -- this is exactly the step that cancels the residual
      \(\tau^{*}(f_{pq}^{*}\varepsilon_p)\) factor produced on the leg-A side;
    \item what remains is the cocycle multiplicativity (C2) of
      \cref{def:scheme_modules_glue} at the triple \((i,p,q)\), conjugated by the
      base-change transports \cref{lem:modules_pullback_basechange_transport}; the
      triple bridges \cref{lem:gr_glueData_bridges} were chosen precisely so that the
      \(\mathrm{pullbackCongr}\) casts on the two sides line up the endpoints on the
      nose.
  \end{enumerate}
  The degenerate pairs \(p = i\) or \(q = i\) need no separate case split: (C2) holds for
  all triples, and where endpoint alignment forces it the self-identity (C1) enters
  through the congruence cast \cref{lem:gr_pullbackCongr_hom_eqToHom}.
\end{proof}

\subsection{The first descent obligation}

\begin{lemma}\leanok
[The candidate-inverse family equalizes the restricted legs (C2 transported)]
  \label{lem:gr_glueChartFamily_equalizes}
  \lean{AlgebraicGeometry.Scheme.Modules.glueChartFamily_equalizes}
  \uses{def:gr_glueChartFamily, def:scheme_modules_glue, lem:gr_glueData_bridges,
    lem:modules_pullback_basechange_transport, lem:gr_overlap_appIso_compat,
    lem:gr_glueChartComponent_self_counit, lem:modules_restrictFunctorIsoPullback_mathlib}
  The family \cref{def:gr_glueChartFamily} equalizes the two restricted descent legs
  \(\iota_i^{*}a\) and \(\iota_i^{*}b\); equivalently, it lands in the restricted
  equalizer \(\iota_i^{*}\mathcal{M}\). This is the first descent obligation, where the
  triple-overlap multiplicativity (C2) of \cref{def:scheme_modules_glue} is consumed.
\end{lemma}
\begin{proof}
  \uses{def:gr_glueChartFamily, lem:gr_glueChartFamily_pullback_map_pi,
    lem:gr_glueChartComponent_leg_compat}
    \leanok
  It suffices to check equality after projecting to each factor \((p,q)\) of the
  restricted overlap product, through the product-preservation comparison of
  \(\iota_i^{*}\) (the same comparison discipline as \cref{lem:gr_glueRestrict_proj_compat}).
  Folding each restricted leg against the \((p,q)\)-projection -- splitting off the
  \(p\)- and \(q\)-th chart projections and cancelling the comparison via
  \cref{lem:gr_glueChartFamily_pullback_map_pi} -- turns the \((p,q)\)-component into the
  equation
  \[
    \mathrm{s}_i^{(p)}\ \circ\ \iota_i^{*}\bigl(a_{(p,q)}\bigr)
      \;=\;
    \mathrm{s}_i^{(q)}\ \circ\ \iota_i^{*}\bigl(b_{(p,q)}\bigr),
  \]
  which is precisely the pair-component obligation
  \cref{lem:gr_glueChartComponent_leg_compat}, where the triple-overlap multiplicativity
  (C2) is consumed. Assembling these componentwise equalities back through the comparison
  gives the equalizing condition.
\end{proof}

\begin{definition}\leanok
[The candidate inverse \(\mathrm{s}_i\)]
  \label{def:gr_glueRestrictionInv}
  \lean{AlgebraicGeometry.Scheme.Modules.glueRestrictionInv}
  \uses{def:gr_glueChartFamily, lem:gr_glueChartFamily_equalizes}
  Because \(\iota_i^{*}\) preserves the descent equalizer, the equalizing family
  \cref{def:gr_glueChartFamily} (equalizing by \cref{lem:gr_glueChartFamily_equalizes})
  lifts uniquely to a morphism
  \[
    \mathrm{s}_i \;:\; \mathcal{M}_i \;\longrightarrow\; \iota_i^{*}\mathcal{M},
  \]
  the candidate inverse of the chart restriction morphism.
\end{definition}

\begin{lemma}\leanok
[Comparison compatibility of the restricted projection]
  \label{lem:gr_glueRestrict_proj_compat}
  \lean{AlgebraicGeometry.Scheme.Modules.glueRestrict_proj_compat}
  \uses{def:gr_glue_equalizer, def:scheme_modules_glue}
  Through the equalizer- and product-preservation comparisons of \(\iota_i^{*}\), the
  restricted equalizer inclusion of \(\iota_i^{*}\mathcal{M}\) into \(\iota_i^{*}P\)
  followed by the \(j\)-th product projection equals the chart pullback
  \(\iota_i^{*}(\pi_j)\) of the \(j\)-th descent-equalizer projection
  (\cref{def:gr_glue_equalizer}).
\end{lemma}
\begin{proof}
  \uses{def:gr_glue_equalizer}
  \leanok
  Pure limit bookkeeping: the equalizer- and product-comparison isomorphisms
  associated to a limit-preserving functor send the restricted inclusion and the
  restricted projection to the pullbacks of the inclusion and the projection
  respectively, and the composite of pullbacks is the pullback of the composite, which
  is \(\iota_i^{*}(\pi_j)\).
\end{proof}

\begin{lemma}\leanok
[The candidate inverse recovers the components]
  \label{lem:gr_glueRestrictionInv_proj}
  \lean{AlgebraicGeometry.Scheme.Modules.glueRestrictionInv_pullback_map_glueProj}
  \uses{def:gr_glueRestrictionInv, lem:gr_glueRestrict_proj_compat,
    def:gr_glueChartComponent}
  The candidate inverse \cref{def:gr_glueRestrictionInv} followed by the restricted
  \(j\)-th projection \(\iota_i^{*}(\pi_j)\) recovers the \(j\)-th component
  \cref{def:gr_glueChartComponent}:
  \[
    \mathrm{s}_i\ \circ\ \iota_i^{*}(\pi_j)
      \;=\; \mathrm{s}_i^{(j)} .
  \]
\end{lemma}
\begin{proof}
  \uses{def:gr_glueRestrictionInv, lem:gr_glueRestrict_proj_compat}
  \leanok
  Unfolding the lift \cref{def:gr_glueRestrictionInv} and applying
  \cref{lem:gr_glueRestrict_proj_compat}, the equalizer- and product-preservation
  comparisons cancel against the lift, leaving the defining \(j\)-th projection of the
  family \cref{def:gr_glueChartFamily}, namely \(\mathrm{s}_i^{(j)}\).
\end{proof}

\begin{lemma}\leanok
[Joint detection of morphisms into the restricted glued sheaf]
  \label{lem:gr_glueRestrict_hom_ext}
  \lean{AlgebraicGeometry.Scheme.Modules.glueRestrict_hom_ext}
  \uses{lem:gr_glueRestrict_proj_compat}
  Two morphisms \(u, v : \mathcal{Z} \to \iota_i^{*}\mathcal{M}\) that agree after every
  restricted projection \(\iota_i^{*}(\pi_j)\) are equal.
\end{lemma}
\begin{proof}
  \uses{lem:gr_glueRestrict_proj_compat}
  \leanok
  The restricted equalizer inclusion is a monomorphism and the restricted product
  projections are jointly monic, both through the preservation comparisons; so it
  suffices to test against \(\iota_i^{*}(\pi_j)\) for every \(j\), which is the
  hypothesis via \cref{lem:gr_glueRestrict_proj_compat}.
\end{proof}

\subsection{The three cocycle obligations}

\begin{lemma}\leanok
[Self-component collapses to the identity (C1 + counit triangle)]
  \label{lem:gr_glueChartComponent_self_counit}
  \lean{AlgebraicGeometry.Scheme.Modules.glueChartComponent_self_counit}
  \uses{def:gr_glueChartComponent, def:scheme_modules_glue,
    lem:gr_restrictIsoPullback_congr, lem:gr_restrictAdjunction_unit_iso,
    lem:gr_restrictIsoPullback_counit, lem:gr_overlapBaseChange_inv_app,
    lem:gr_pullbackCongr_hom_eqToHom}
  The \((i,i)\)-component \cref{def:gr_glueChartComponent}, transposed back along the
  chart immersion \(\iota_i\) (post-composed with the geometric counit), is the
  identity of \(\mathcal{M}_i\):
  \[
    \mathrm{s}_i^{(i)}\ \circ\ \varepsilon^{\mathrm{geo}}_{\mathcal{M}_i}
      \;=\; \mathrm{id}_{\mathcal{M}_i}.
  \]
  This is the second descent obligation, consuming the self-identity (C1).
\end{lemma}
\begin{proof}
  \uses{def:gr_glueChartComponent, def:scheme_modules_glue,
    lem:gr_restrictIsoPullback_congr, lem:gr_restrictAdjunction_unit_iso,
    lem:gr_restrictIsoPullback_counit, lem:gr_overlapBaseChange_inv_app,
    lem:gr_pullbackCongr_hom_eqToHom}
    \leanok
  Since \(f_{ii}\) factors as \(t_{ii}\circ f_{ii}\) (because \(t_{ii} = \mathrm{id}\)),
  (C1) of \cref{def:scheme_modules_glue} says the transition \(g_{ii}\) is the
  canonical congruence cast; through \cref{lem:gr_pullbackCongr_hom_eqToHom} and
  \cref{lem:gr_restrictIsoPullback_congr} the three middle factors of the component
  collapse to a single site-level congruence cast. Rewriting the geometric unit by
  \cref{lem:gr_restrictAdjunction_unit_iso} and the geometric counit by
  \cref{lem:gr_restrictIsoPullback_counit} turns the whole composite into the
  site-level unit, the congruence cast, the section form of \(\beta_{ii}^{-1}\)
  (\cref{lem:gr_overlapBaseChange_inv_app}), and the site-level counit. On sections
  over an open this is a cycle of presheaf restrictions along opens identities whose
  composite open-inclusion is the identity; folding the restrictions and using
  proof-irrelevance of the open inclusions gives the identity.
\end{proof}

\begin{lemma}\leanok
[Transpose of the overlap base change (mate core)]
  \label{lem:gr_glueOverlapFactor_transpose}
  \lean{AlgebraicGeometry.Scheme.Modules.glueOverlapFactor_transpose}
  \uses{def:gr_glueOverlapFactorIso, lem:glueOverlapBaseChangeIso,
    lem:gr_overlapBaseChange_hom_app, def:scheme_modules_glue}
  The adjoint transpose along \(\iota_i\) of the pullback-form overlap base change
  \(\gamma_{ij}\) (\cref{def:gr_glueOverlapFactorIso}) is the canonical four-functor
  comparison of the glue square:
  \[
    \widehat{\gamma_{ij}}
      \;=\;
    (\iota_j)_{*}\bigl(\eta_{t_{ij}\circ f_{ji}}\bigr)
      \ \circ\ \mathrm{pushComp}
      \ \circ\ \mathrm{pushCongr}
      \ \circ\ \mathrm{pushComp}^{-1},
  \]
  the unit along \(t_{ij}\circ f_{ji}\) pushed forward along \(\iota_j\), regrouped by
  the pushforward-composition comparison, reindexed by the pushforward-congruence along
  the glue condition \((t_{ij}\circ f_{ji})\circ \iota_j = f_{ij}\circ \iota_i\), and
  ungrouped by the inverse pushforward-composition. The identity is free of the
  cocycle hypotheses: it is pure site-level content.
\end{lemma}
\begin{proof}
  \uses{def:gr_glueOverlapFactorIso, lem:glueOverlapBaseChangeIso,
    lem:gr_overlapBaseChange_hom_app, def:scheme_modules_glue,
    lem:modules_restrictFunctorIsoPullback_mathlib}
    \leanok
  Expand the adjoint transpose of \(\gamma_{ij}\): the left-adjoint comparison
  isomorphisms (\cref{lem:modules_restrictFunctorIsoPullback_mathlib}) conjugating the
  chart factor cancel against the unit/counit, leaving the bare site-level overlap base
  change \(\beta_{ij}\) (\cref{lem:glueOverlapBaseChangeIso}). The unit bridge at
  \(t_{ij}\circ f_{ji}\) and the three naturality squares of the pushforward-composition
  and pushforward-congruence comparisons rewrite the transpose into the stated
  four-factor composite up to whiskering and reassociation. The residual site-level
  core is checked on sections over an open: by \cref{lem:gr_overlapBaseChange_hom_app}
  the section form of \(\beta_{ij}\) is a presheaf restriction along an opens identity,
  and the remaining composite is a cycle of such restrictions whose total open
  inclusion is the identity, folded by proof-irrelevance exactly as in the closing step
  of \cref{lem:gr_glueChartComponent_self_counit}.
\end{proof}

\begin{lemma}\leanok
[Mate recognition of the factor transpose]
  \label{lem:gr_glueOverlapFactor_mate}
  \lean{AlgebraicGeometry.Scheme.Modules.glueOverlapFactor_mate}
  \uses{lem:gr_glueOverlapFactor_transpose, def:gr_glueOverlapFactorIso,
    def:scheme_modules_glue}
  For any morphism \(m : \mathcal{W} \to (\iota_j)_{*}\mathcal{M}_j\) on the glued
  scheme, the conjugated transport of the inverse transpose of \(m\) along the glue
  square equals the inverse transpose, along \(f_{ij}\), of \(\iota_i^{*}(m)\) followed
  by \(\gamma_{ij}\) (\cref{def:gr_glueOverlapFactorIso}). In other words, the factor
  transpose \cref{lem:gr_glueOverlapFactor_transpose} is the mate of \(\gamma_{ij}\)
  through the pseudofunctor structure of pullback.
\end{lemma}
\begin{proof}
  \uses{lem:gr_glueOverlapFactor_transpose, def:gr_glueOverlapFactorIso,
    def:scheme_modules_glue}
    \leanok
  Both sides are computed by transposing twice. By injectivity of the two transpose
  bijections it suffices to compare the double transposes. Naturality of the transpose
  in its left argument peels \(m\) off the right-hand side; the conjugate of the
  pullback-composition isomorphism along the glue condition is the inverse
  pushforward-composition isomorphism; and folding the unit pair and reindexing along
  the glue condition of \cref{def:scheme_modules_glue} reduces the identity to the
  factor transpose \cref{lem:gr_glueOverlapFactor_transpose}.
\end{proof}

\begin{lemma}\leanok
[Overlap compatibility of the restriction morphisms]
  \label{lem:gr_glueRestriction_overlap_compat}
  \lean{AlgebraicGeometry.Scheme.Modules.glueRestriction_overlap_compat}
  \uses{lem:glueRestrictionHom, def:scheme_modules_glue, def:gr_glue_equalizer}
  The \((i,j)\)-component of the descent-equalizer condition, transposed to the
  pullback level, states that over the overlap \(V_{ij}\) the two restriction
  morphisms \(\mathrm{r}_i\) and \(\mathrm{r}_j\) (\cref{lem:glueRestrictionHom}) differ
  exactly by the transition isomorphism \(g_{ij}\) (through the pseudofunctor
  congruence casts).
\end{lemma}
\begin{proof}
  \uses{lem:glueRestrictionHom, def:scheme_modules_glue, def:gr_glue_equalizer}
  \leanok
  The glued sheaf is the descent equalizer, so its \(i\)- and \(j\)-projections satisfy
  the overlap-agreement condition of \cref{def:gr_glue_equalizer}. Transposing that
  condition along \(\iota_i\) and \(\iota_j\) and folding the transposes back through
  the adjunction bijection turns it into the stated identity between the overlap
  pullbacks of \(\mathrm{r}_i\) and \(\mathrm{r}_j\), conjugated by the glue-condition
  casts of \cref{def:scheme_modules_glue} and corrected by \(g_{ij}\).
\end{proof}

\begin{lemma}\leanok
[Pair condition transposed: restriction composed with the component]
  \label{lem:gr_glueRestrictionHom_glueChartComponent}
  \lean{AlgebraicGeometry.Scheme.Modules.glueRestrictionHom_glueChartComponent}
  \uses{lem:gr_glueRestriction_overlap_compat, lem:gr_glueOverlapFactor_mate,
    def:gr_glueChartComponent, lem:glueRestrictionHom}
  The chart restriction morphism \(\mathrm{r}_i\) (\cref{lem:glueRestrictionHom})
  followed by the \(j\)-th candidate-inverse component \cref{def:gr_glueChartComponent}
  equals the restricted \(j\)-th projection:
  \[
    \mathrm{r}_i\ \circ\ \mathrm{s}_i^{(j)}
      \;=\; \iota_i^{*}(\pi_j).
  \]
  This is the third descent obligation; together with
  \cref{lem:gr_glueRestrictionInv_proj} and \cref{lem:gr_glueRestrict_hom_ext} it gives
  \(\mathrm{r}_i \circ \mathrm{s}_i = \mathrm{id}\).
\end{lemma}
\begin{proof}
  \uses{lem:gr_glueRestriction_overlap_compat, lem:gr_glueOverlapFactor_mate,
    def:gr_glueChartComponent, lem:glueRestrictionHom}
    \leanok
  Unfold the component \cref{def:gr_glueChartComponent}. By the overlap compatibility
  \cref{lem:gr_glueRestriction_overlap_compat} the overlap pullback of \(\mathrm{r}_i\)
  transported across \(g_{ij}\) is the inverse mate transpose of the restricted
  projection \(\iota_i^{*}(\pi_j)\) through \(\gamma_{ij}\)
  (\cref{lem:gr_glueOverlapFactor_mate}). Substituting this into the unit-naturality
  square of \(\mathrm{r}_i\) along \(f_{ij}\) and folding the resulting transpose
  cancels the factor isomorphism, leaving exactly \(\iota_i^{*}(\pi_j)\).
\end{proof}

\subsection{Joint faithfulness of chart restrictions}

\begin{lemma}\leanok
[Chart restrictions are jointly faithful]
  \label{lem:gr_pullback_jointly_faithful}
  \lean{AlgebraicGeometry.Scheme.Modules.pullback_map_jointly_faithful}
  \uses{lem:modules_restrictFunctorIsoPullback_mathlib, lem:eq_of_locally_eq_mathlib,
    def:scheme_modules_glue}
  Two morphisms \(u, v : \mathcal{W} \to \mathcal{W}'\) of sheaves of modules on the
  glued scheme that agree after pullback to every chart, \(\iota_i^{*}u = \iota_i^{*}v\)
  for all \(i\), are equal. This realises the restriction-uniqueness core of the
  characterisation of the glued sheaf up to unique isomorphism: a morphism out of the
  glued sheaf is determined chart-locally.
\end{lemma}
\begin{proof}
  \uses{lem:modules_restrictFunctorIsoPullback_mathlib, lem:eq_of_locally_eq_mathlib}
  \leanok
  Transfer the chart agreement to the site-level restriction functor by conjugating
  with the comparison \cref{lem:modules_restrictFunctorIsoPullback_mathlib}, whose
  sections are concrete: \(\Gamma(\iota_i^{\mathrm{res}}\mathcal{W}, V) =
  \Gamma(\mathcal{W}, \iota_i(V))\). Then test \(u\) and \(v\) on sections over any
  open \(O\) of the glued scheme using separatedness of the target sheaf
  (\cref{lem:eq_of_locally_eq_mathlib}) over the cover \(\{\iota_i(\iota_i^{-1}O)\}\)
  of \(O\) — which is a cover by joint surjectivity of the chart immersions
  (\cref{def:scheme_modules_glue}). On each cover piece the agreement is the chart
  identity moved inside the sections by naturality of \(u, v\) on the presheaf level.
\end{proof}

\begin{definition}

\leanok
[The restriction isomorphism of the glued sheaf]
  \label{def:glueRestrictionIso}
  \lean{AlgebraicGeometry.Scheme.Modules.glueRestrictionIso}
  % NOTE: this is the 1:1 home for the Lean decl `glueRestrictionIso`
  % (AlgebraicJacobian/Picard/GlueDescent.lean). Reconciliation complete as of iter-080:
  % `def:gr_modules_glueRestrictionIso` in Picard_GrassmannianQuot.tex now carries NO
  % `\lean{}` pin and is a cross-reference node only, so the pin is no longer duplicated.
  \uses{thm:isIso_glueRestrictionHom}
  Let \(\mathcal{M}\) be the glued sheaf of \cref{def:scheme_modules_glue}, with the
  transition family \((g_{ij})\) satisfying (C1)/(C2). For each chart index \(i\) the
  restriction morphism \(\mathrm{r}_i\) (\cref{lem:glueRestrictionHom}) is promoted to
  the \emph{restriction isomorphism}
  \[
    \rho_i \;:\; \iota_i^{*}\,\mathcal{M} \;\xrightarrow{\ \sim\ }\; \mathcal{M}_i,
  \]
  the canonical identification of the chart restriction of the glued sheaf with the
  \(i\)-th input sheaf, packaged from \cref{thm:isIso_glueRestrictionHom}. Its
  underlying morphism is the adjoint transpose of the \(i\)-th descent-equalizer
  projection, and over each overlap \(V_{ij}\) the two restrictions \(\rho_i, \rho_j\)
  differ exactly by the transition isomorphism \(g_{ij}\).
\end{definition}
